\usepackage{fontspec}
\usepackage{polyglossia}
\setdefaultlanguage{russian}
\setotherlanguage{english}
\setmainfont{Liberation Serif}
\setsansfont{Liberation Sans}
\setmonofont{Liberation Mono}

\newfontfamily\cyrillicfontsf{Liberation Sans}

\usepackage[
  left=30mm,
  right=15mm,
  top=20mm,
  bottom=20mm
]{geometry}

% Полуторный межстрочный интервал
\usepackage{setspace}
\onehalfspacing

% Красная строка (абзацный отступ) 1.25 см
\usepackage{indentfirst}
\setlength{\parindent}{1.25cm}

\usepackage{titlesec}
\titleformat{\section}{\normalfont\Large\bfseries\centering}{\thesection}{1em}{}
\titleformat{\subsection}{\normalfont\large\bfseries\raggedright}{\thesubsection}{1em}{}
\titleformat{\subsubsection}{\normalfont\normalsize\bfseries\raggedright}{\thesubsubsection}{1em}{}
\let\oldsection\section
\renewcommand{\section}{\clearpage\oldsection}

% Отмена переносов в заголовках
\usepackage{ragged2e}

% Настройка нумерации страниц (внизу по центру)
\usepackage{fancyhdr}
\pagestyle{fancy}
\fancyhf{}
\cfoot{\thepage}
\renewcommand{\headrulewidth}{0pt}
% Настройки для автоматического переноса строк в блоках кода (listings)
\usepackage{listings}
\usepackage{xcolor}

\lstset{
  breaklines=true,
  breakatwhitespace=false,
  basicstyle=\ttfamily\small,
  frame=single,
  xleftmargin=0pt,
  xrightmargin=0pt
}
\usepackage{fvextra}
\DefineVerbatimEnvironment{Highlighting}{Verbatim}{
  breaksymbol=,
  breaklines=true,
  commandchars=\\\{\}
}
